\section{Access}\label{sec:access}

The Access chapter is about what an instrument can distinguish.  An
observation signature \(\obsSig\) (\Cref{def:access-quotient}) is a
declared map from states to observation symbols; its access quotient
\(\accessQ\) identifies exactly those states that no declared
observation separates.  The laws of this chapter say how such
signatures determine quotients, when those quotients are adequate for
accepted claims, when a new distinction requires new apparatus data
rather than a free relabeling, and when structure remains hidden behind
the current quotient \(\currQ\) although a predictive quotient
\(\predM\) separates it.  In the running example of \Cref{sec:intro},
\(q\) is the access quotient of the signature that reads each state's
class \(a\), \(b\) or \(c\); the states \(1\) and \(2\) are
indistinguishable to it, which is why \(F\) cannot be read off \(q\).

Quotientality (\Fref{F10}) says that an observation signature
determines a canonical access quotient, through which every observable
constant on the signature factors.  Adequacy / No-Overread (\Fref{F11})
says that accepted claims must factor through the available access.
No-Free-Distinction (\Fref{F12}) says that a readout splits an access
fiber exactly when it does not factor through the access quotient, that
such a readout cannot be accepted as exact under the no-overread rule,
and that the pairing \(\langle\accessQ,d\rangle\) is the coarsest
refinement carrying it.  Layer Multiplicity (\Fref{F24}) closes the
chapter: a role readout that the access quotient does not determine is
carried by a strict refinement of that quotient, an additional layer.

Between these, the Hiddenness Normal Form (\Fref{F13a}) and the CSL
Formed-Closure Law (\Fref{F13b}) compare current and predictive
observation.  Hiddenness is the situation in which states agree under
\(\currQ\) but differ under \(\predM\).  \Fref{F13a} is conditional on a
source hypothesis motivated by \citet{Tsiokos2026PvNP}; \Fref{F13b} is
a direct theorem under an explicit determining hypothesis, motivated by
\citet{Tsiokos2026CSL}.  No result of this paper assumes the
P~\(\neq\)~NP claim of the former.  On a first reading these two laws
can be skipped; the rest of the chapter does not use them.

\subsection{Quotientality (Layer-Sameness via Quotient) [\texorpdfstring{\Fref{F10}}{F10}]}\label{sec:access:quotientality}

\Cref{def:obs-signature} restates \Cref{def:access-quotient} for a closure
carrier, names the quotient set \(Q_I\), and adds the calibration-family case.

\begin{definition}[Observation signature]\label{def:obs-signature}
Let \(H\) be a closure carrier.  An observation signature on \(H\)
is a declared map
\[
  \obsSig:H\to\Sigma
\]
from states to observation symbols.  Its access kernel is the
equivalence relation
\[
  h\sim_I h'
  \quad\Longleftrightarrow\quad
  \obsSig(h)=\obsSig(h'),
\]
and its access quotient is the quotient map
\[
  \accessQ:H\to Q_I:=H/{\sim_I}.
\]
For a family \((\sigma_j)_{j\in J}\) of observation signatures (a calibration
family), the joint access kernel is \(h\sim h'\iff\sigma_j(h)=\sigma_j(h')\) for
every \(j\), the kernel of the bundled signature \(h\mapsto(\sigma_j(h))_j\), and
the same construction is applied to it.
\end{definition}

\begin{theorem}[Quotientality]\label{thm:F10}
Let \(\obsSig:H\to\Sigma\) be an observation signature with access
quotient \(\accessQ:H\to Q_I\).  Every observable \(o:H\to R\) that
is constant on access classes factors uniquely through
\(\accessQ\):
\[
  o=\overline o\circ \accessQ
\]
for a unique \(\overline o:Q_I\to R\).  Moreover, \(Q_I\) is the
coarsest quotient carrying the declared observations: if
\(\pi:H\to Q'\) is any surjective map on whose fibers \(\obsSig\) is constant,
then \(\accessQ\) factors through \(\pi\).
\end{theorem}
\begin{proof}
Define \(\overline o([h]_I)=o(h)\).  This is well defined because
\([h]_I=[h']_I\) means \(h\sim_I h'\), and the hypothesis on \(o\)
gives \(o(h)=o(h')\).  The equation
\(o=\overline o\circ\accessQ\) follows from the definition on
representatives.  If \(\overline o'\) is another such map, then for
every access class \([h]_I\) we have
\[
  \overline o'([h]_I)
  =
  \overline o'(\accessQ(h))
  =
  o(h)
  =
  \overline o([h]_I),
\]
so \(\overline o'=\overline o\).

For coarsestness, suppose \(\pi:H\to Q'\) is a surjective map on whose fibers
\(\obsSig\) is constant: \(\pi(h)=\pi(h')\) implies
\(\obsSig(h)=\obsSig(h')\).  Since \(\pi\) is surjective, define
\(u:Q'\to Q_I\) by \(u(\pi(h))=\accessQ(h)\).  The same invisibility
hypothesis makes this assignment independent of the representative.
Hence \(\accessQ=u\circ\pi\).  Thus any quotient already carrying the
declared observations refines the canonical access quotient.
\end{proof}

\paragraph{Nonclaims.} The
theorem does not produce the coarsest quotient possible for all
observers; it produces the coarsest quotient carrying the given
signature.  It also does not create new observable distinctions and
does not instantiate any particular carrier theory.

\subsection*{Layer instantiations}\label{layer-inst:F10}

\begingroup\small
\begin{description}
\item[\emph{Myhill--Nerode minimal {DFA} (formal language theory).}] \textit{Special case.}
For a regular language \(L\subseteq\Sigma^*\), take \(H=\Sigma^*\) and
\(\obsSig(w):=w^{-1}L=\{z:wz\in L\}\), valued in the power set of \(\Sigma^*\) (here \(\Sigma\) is the alphabet, not the symbol set of \Cref{def:obs-signature}). Its access kernel is the Nerode
relation \(\equiv_L\), defined by
\(w\equiv_L w'\Leftrightarrow\forall z.\,(wz\in L\Leftrightarrow w'z\in L)\),
so \(Q_I=\Sigma^*/\!\!\equiv_L\) is the state set of the minimal DFA
recognising \(L\). \Fref{F10} then says that every observable constant on
\(\equiv_L\) factors uniquely through \(Q_I\), and that every surjection with
\(\obsSig\)-constant fibers refines \(\equiv_L\). That \(\equiv_L\) is also the
coarsest right-congruence making \(\mathbf 1_L\) class-constant is the
Myhill--Nerode theorem, not part of \Fref{F10}
\citep{Nerode1958LinearAutomatonTransformations}.

\item[\emph{Kalman observability canonical decomposition (control
engineering).}] \textit{Special case.}
Take \(H=\mathbb R^n\) and \(\obsSig(x):=(CA^kx)_{k\ge0}\); its access kernel
is \(x-x'\in N\). For a linear time-invariant system \(\dot x=Ax+Bu,\ y=Cx\), the
unobservable subspace
\(N=\bigcap_{k\ge 0}\ker(CA^k)\) is the kernel of the observability
matrix.  Two state vectors agree under \(\mathbb R^n/N\) iff they
produce identical output trajectories for every admissible input,
and the Kalman canonical form factors every input--output
observable uniquely through \(\mathbb R^n/N\); every surjection of \(\mathbb R^n\) on whose fibres the output trajectory is constant refines \(\mathbb R^n/N\) (the quotient map \(\mathbb R^n\to\mathbb R^n/N\) factors through it), which is the coarsest-quotient clause of \Fref{F10}
\citep{Kalman1963MathematicalDescription}.

\item[\emph{Bernstein 3{NF} synthesis (database engineering).}] \textit{Analogy.}
Given a relation schema with attributes \(U\) and a
functional-dependency set \(F\), the 3NF decomposition produced by Bernstein
synthesis from a chosen minimal cover of \(F\) is a
dependency-preserving 3NF decomposition, which depends on the chosen
cover: every FD-checker for the cover factors through the projection
onto a single component schema, and coalescing components keeps
every declared FD checkable but can introduce a 3NF violation
\citep{Bernstein1976Synthesizing3NF}.

\end{description}
\endgroup


\subsection{Adequacy / No-Overread [\texorpdfstring{\Fref{F11}}{F11}]}\label{sec:access:adequacy}

Fix the same closure carrier \(H\) and declared observable family as
in the previous subsection, and let \(\accessQ:H\to Q_I\) be the
access quotient constructed from that family.  A claim is \emph{accepted as an exact current claim} when it passes
the finite audited status discipline of \FoundIII\ as an exact claim
about the current layer \citep{Tsiokos2026FoundationsIII}.  By the
no-overread rule (NO) of \Cref{sec:apparatus}, it may then use only the
distinctions available to the current access; this is the premise the
proof uses.  The theorem records what the discipline guarantees: an
accepted exact claim is constant on the fibers of \(\accessQ\), and
hence factors through it.

\begin{theorem}[Adequacy / No-Overread]\label{thm:F11}
Let \(F:H\to W\) be a map (in the intended reading, a current-layer claim over a
formed access package), and let
\(\mathcal O(\accessQ,F):=\{(h,h'):\accessQ(h)=\accessQ(h'),\ F(h)\ne F(h')\}\)
be its overread obstruction.
\begin{enumerate}[label=(\alph*)]
\item For every \(F\),
\[
  F\ \text{constant on the fibers of}\ \accessQ
  \iff \mathcal O(\accessQ,F)=\emptyset
  \iff F=\overline F\circ\accessQ\ \text{for a unique}\ \overline F:Q_I\to W .
\]
\item Under the no-overread rule (NO) of \Cref{sec:apparatus}, if \(F\) is
  accepted as an exact current claim, then it is constant on the fibers of
  \(\accessQ\), and there is a unique map \(\overline F:Q_I\to W\) such that
  \(F=\overline F\circ\accessQ\).
\item Hence the no-overread rule (NO) holds if and only if every accepted
  exact current claim factors uniquely through
  \(\accessQ\).
\end{enumerate}
The first conclusion of (b) is (NO) applied to \(F\); the uniqueness uses the
universal property of \(\accessQ\) (\Cref{thm:F10}).
\end{theorem}
\begin{proof}
For (b), (NO) of \Cref{sec:apparatus} applied to the accepted claim \(F\) states exactly that \(F(h)=F(h')\) whenever \(\accessQ(h)=\accessQ(h')\).

By Quotientality, every map constant on access classes factors
uniquely through \(\accessQ\).  Applying \Cref{thm:F10} to \(F\)
gives a unique \(\overline F:Q_I\to W\) with
\(F=\overline F\circ\accessQ\).  Thus an accepted current claim has
no overread defect. For the equivalences, a pair in \(\mathcal O(\accessQ,F)\)
is exactly a failure of constancy on a fiber, and Quotientality, with
surjectivity of \(\accessQ\), gives the unique factorization of a map constant
on fibers; conversely a factored map is constant on fibers. Applying these
equivalences to each accepted exact current claim gives the final
statement.
\end{proof}

\paragraph{Nonclaims.} The theorem does not forbid recording distinctions beyond
current access; such distinctions may belong to the predictive
quotient \(\predM\), to memory, or to a later access extension.  It
also does not determine the audit data, bridge strengths, thresholds,
or residual budgets required for acceptance.

\subsection*{Layer instantiations}\label{layer-inst:F11}

\begingroup\small
\begin{description}
\item[\emph{Rao--Blackwell theorem (statistical decision theory).}] \textit{Analogy.}
For a sufficient statistic \(T:X\to\mathcal T\) on a parametric
family \(\{P_\theta\}\), the Rao--Blackwell construction
\(\delta_{\mathrm{RB}}(x)=\mathbb E[\delta(X)\mid T(X)=T(x)]\) takes any
decision rule \(\delta\) and produces a \(\sigma(T)\)-measurable
rule whose risk under a convex loss is at most that of \(\delta\)
at every \(\theta\), and strictly smaller at some \(\theta\) under a
strictly convex loss unless \(\delta\) is already a function of \(T\)
almost surely.  Hence under a strictly convex loss any admissible
decision rule with finite risk factors as
\(\delta=\bar\delta\circ T\) a.s.; a rule that uses within-fibre
detail of \(X\) beyond \(T(X)\) is the overread failure
\citep{Blackwell1947ConditionalExpectation}.

\end{description}
\endgroup


\subsection{No-Free-Distinction [\texorpdfstring{\Fref{F12}}{F12}]}\label{sec:access:no-free-distinction}

Keep the formed access quotient \(\accessQ:H\to Q_I\) from
Quotientality, together with the no-overread discipline from
Adequacy.  A proposed new current distinction is a readout
\(d:H\to D\).  Its split-pair obstruction at fixed access is
\[
  \splitObs_{\mathrm{dist}}(\accessQ,d)
  =
  \{(h,h')\mid \accessQ(h)=\accessQ(h')\ \text{and}\ d(h)\ne d(h')\}.
\]

Call \(d\) \emph{current-visible} when \(\accessQ(h)=\accessQ(h')\) implies
\(d(h)=d(h')\). The \emph{financed refinement} of \(d\), the refinement that adjoins \(d\)
as declared apparatus data, is
\(\langle\accessQ,d\rangle:H\to Q_I\times D\).

\begin{theorem}[No-Free-Distinction]\label{thm:F12}
Let \(d:H\to D\) be a readout (the access quotient \(\accessQ\) is surjective by construction). Then
\(d\) is current-visible if and only if
\[
  \splitObs_{\mathrm{dist}}(\accessQ,d)=\varnothing,
  \qquad\text{equivalently}\qquad
  d=\overline d\circ\accessQ
\]
for some \(\overline d:Q_I\to D\); in particular a readout obtained by
applying a map to the access quotient cannot split an access fiber. The
financed refinement \(\langle\accessQ,d\rangle\) retains \(\accessQ\) and
carries \(d\) (both factor through it), it strictly refines \(\accessQ\) when
the obstruction is nonempty (it separates two histories that \(\accessQ\)
identifies, so it does not factor through \(\accessQ\)), and every map \(p\) through which both
\(\accessQ\) and \(d\) factor also has \(\langle\accessQ,d\rangle\) factoring
through it. If the obstruction is nonempty and the no-overread rule (NO) of
\Cref{sec:apparatus} holds, \(d\) is not accepted as an exact current claim.
\end{theorem}
\begin{proof}
By definition, \(d\) is current-visible exactly when no pair identified by
\(\accessQ\) is separated by \(d\), that is, when the displayed obstruction is
empty. By \Cref{thm:F10}, this is equivalent to the existence of a
factorization \(d=\overline d\circ\accessQ\). A readout \(g\circ\accessQ\)
depends only on the access class, so it is current-visible.

The projections of \(Q_I\times D\) give
\(\accessQ=\operatorname{pr}_1\circ\langle\accessQ,d\rangle\) and
\(d=\operatorname{pr}_2\circ\langle\accessQ,d\rangle\). If
\((h,h')\) lies in the obstruction, then \(\accessQ(h)=\accessQ(h')\) while
\(\langle\accessQ,d\rangle(h)\ne\langle\accessQ,d\rangle(h')\), so the
refinement is strict. If \(\accessQ=a\circ p\) and \(d=b\circ p\), then
\(\langle\accessQ,d\rangle=\langle a,b\rangle\circ p\).

Finally, if the obstruction is nonempty, \(d\) separates two histories that
the current access quotient identifies, so \(d\) is not constant on the fibers
of \(\accessQ\). By (NO), as in \Cref{thm:F11}, an exact current claim is
constant on those fibers; hence \(d\) is not accepted as exact.
\end{proof}

A distinction that the access quotient does not support can still be used
with declared apparatus data: memory, a bridge, a budgeted residual, a
scope restriction, or stronger calibration. Among quotient refinements
through which both \(\accessQ\) and \(d\) factor,
\(\langle\accessQ,d\rangle\) is coarsest, by the factorization proved above.

\paragraph{Nonclaims.} The theorem is not anti-refinement.  Calibration families may enrich
the apparatus, and strict extensions may add real distinctions.  The
claim is only that the enrichment must be declared and statused, not
smuggled in as a free consequence of the old access quotient.  Nor
does this theorem classify every possible source of new data; later
axes treat memory, bridges, residual budgets, scope changes, and
predictive refinements in their own normal forms.

\subsection*{Layer instantiations}\label{layer-inst:F12}

\begingroup\small
\begin{description}
\item[\emph{Landauer's principle (thermodynamics of computation).}] \textit{Analogy.}
Resetting a one-bit register whose prior content is an unknown,
unbiased bit to a fixed value, by a process that exchanges heat with a
thermal reservoir at temperature \(T\) and starts uncorrelated with it,
dissipates at least \(k_B T\ln 2\) of heat into the reservoir; in
general the bound is \(k_BT\) times the decrease of the register's
entropy, measured in nats, so the lower bound is smaller for a biased prior bit
\citep{Landauer1961IrreversibilityHeatGeneration,ReebWolf2014ImprovedLandauer}.
The analogy is that a change in the distinctions a register carries is
not free: it is paid for by declared apparatus, here heat exported to
the reservoir.

\end{description}
\endgroup



\subsection{Hiddenness Normal Form (conditional) [\texorpdfstring{\Fref{F13a}}{F13a}]}\label{sec:access:hiddenness}

\begin{definition}[Hiddenness interface]\label{def:hiddenness-exposure}
A \emph{hiddenness interface} \(I\) consists of a closure carrier
\(\mathcal H_I\), a current quotient
\(\currQ_I:\mathcal H_I\to Q_I^{\mathrm{cur}}\), a predictive quotient
\(\predM_I:\mathcal H_I\to M_I\) refining \(\currQ_I\), an upstream
readout \(u_I:\mathcal H_I\to U_I\) of the determining state, and three
declared propositions
\[
  \operatorname{determiningStateTyped}(I),\quad
  \operatorname{minimalRelevantUpstream}(I),\quad
  \operatorname{lawfulCurrentExposure}(I).
\]
From these data we define \(\operatorname{QuotientExposed}(I)\),
\(\operatorname{PredictiveCollapse}(I)\), \(\operatorname{PredictiveSurplus}(I)\)
and \(\operatorname{HiddenFromCurrent}(I)\) below.
\Fref{F13b} defines a
distinct pair-level predicate using a readout \(d\) and proves an equivalence
under an explicit determining hypothesis. \(\operatorname{QuotientExposed}\),
\(\operatorname{PredictiveCollapse}\), \(\operatorname{PredictiveSurplus}\) and
\(\operatorname{HiddenFromCurrent}\) are defined from the maps and the declared
predicates:
\begin{align*}
  \operatorname{QuotientExposed}(I)
  &:\Longleftrightarrow \exists f,\ u_I=f\circ\currQ_I,\\
  \operatorname{PredictiveCollapse}(I)
  &:\Longleftrightarrow \exists g,\ \predM_I=g\circ\currQ_I,\\
  \operatorname{PredictiveSurplus}(I)
  &:\Longleftrightarrow \exists h,h',\ \currQ_I(h)=\currQ_I(h')
    \wedge \predM_I(h)\ne\predM_I(h'),\\
  \operatorname{HiddenFromCurrent}(I)
  &:\Longleftrightarrow \neg\operatorname{lawfulCurrentExposure}(I).
\end{align*}
The hiddenness normal form is the
five-conjunct property
\begin{align*}
  \operatorname{HiddennessNormalForm}(I)
  :={}&
  \operatorname{determiningStateTyped}(I)
  \wedge
  \operatorname{minimalRelevantUpstream}(I)\\
  &{}\wedge
  \bigl(\operatorname{lawfulCurrentExposure}(I)
    \leftrightarrow
    \operatorname{QuotientExposed}(I)\bigr)\\
  &{}\wedge
  \bigl(\operatorname{PredictiveCollapse}(I)
    \leftrightarrow
    \operatorname{lawfulCurrentExposure}(I)\bigr)\\
  &{}\wedge
  \bigl(\operatorname{PredictiveSurplus}(I)
    \leftrightarrow
    \operatorname{HiddenFromCurrent}(I)\bigr).
\end{align*}
\end{definition}

\paragraph{Reading the statement.}
The current quotient \(\currQ_I\) records what is observable now, and
the predictive quotient \(\predM_I\) records what matters for
prediction.  \(\operatorname{PredictiveCollapse}(I)\) says that
\(\predM_I\) factors through \(\currQ_I\), so that prediction needs
nothing beyond the current class; \(\operatorname{PredictiveSurplus}(I)\)
says that two histories with the same current class have different
predictive classes; and \(\operatorname{QuotientExposed}(I)\) says that
the whole upstream readout factors through \(\currQ_I\).  The predicates
\(\operatorname{determiningStateTyped}(I)\) (the determining state has
been given a type), \(\operatorname{minimalRelevantUpstream}(I)\) (the
relevant upstream data are minimal) and
\(\operatorname{lawfulCurrentExposure}(I)\) (the current quotient is a
lawful exposure of the determining state) are declared data of the
interface, not reduced to the quotients.

The theorem has two parts.  Part (i) assumes the source hypothesis
\((\mathsf S)\): \(\currQ_I\) is surjective and the first four conjuncts
of the normal form hold.  It proves the fifth conjunct and the chain of
equivalences linking exposure, collapse and the absence of surplus.  In
the intended application, \(I\) would be an interface extracted from
the saturated-SAT construction of \citet{Tsiokos2026PvNP}, and
\((\mathsf S)\) a property of that construction; that paper constructs
no such interface, and \((\mathsf S)\) is not verified for it here.
Hypothesis \((\mathsf S)\) is satisfiable, also together with a
predictive surplus: take two histories, a one-point current quotient,
the identity as predictive quotient and as upstream readout,
\(\operatorname{lawfulCurrentExposure}(I)\) false and the other two
declared predicates true.  Part (ii) assumes nothing beyond the
interface: the equivalence of surplus with non-collapse holds exactly
when \(\currQ_I\) is surjective or \(M_I\) is nonempty.  The refinement
of \(\currQ_I\) by \(\predM_I\) required in
\Cref{def:hiddenness-exposure} is not used by the theorem; it serves
only the remark, in the nonclaims below, that
\(\langle\currQ_I,\predM_I\rangle\) has the fibers of \(\predM_I\).

\begin{theorem}[Hiddenness Normal Form]\label{thm:F13a}
Let \(I\) be a hiddenness interface.
\begin{enumerate}[label=(\roman*)]
\item Assume the source hypothesis
  \((\mathsf S)\): \(\currQ_I\) is surjective, and
  \begin{align*}
    &\operatorname{determiningStateTyped}(I),\\
    &\operatorname{minimalRelevantUpstream}(I),\\
    &\operatorname{lawfulCurrentExposure}(I)
      \iff \operatorname{QuotientExposed}(I),\\
    &\operatorname{PredictiveCollapse}(I)
      \iff \operatorname{lawfulCurrentExposure}(I).
  \end{align*}
  Then
  \[
    \operatorname{PredictiveSurplus}(I)
      \iff \operatorname{HiddenFromCurrent}(I),
  \]
  so \(I\) satisfies \(\operatorname{HiddennessNormalForm}(I)\), and
  \[
    \operatorname{QuotientExposed}(I)\iff
    \operatorname{PredictiveCollapse}(I)\iff
    \neg\operatorname{PredictiveSurplus}(I):
  \]
  the upstream readout descends through the current quotient exactly
  when prediction does.
\item Without \((\mathsf S)\), \(I\) satisfies
  \(\operatorname{PredictiveSurplus}(I)\iff\neg\operatorname{PredictiveCollapse}(I)\)
  if and only if \(\currQ_I\) is surjective or \(M_I\) is nonempty, in
  particular whenever \(\currQ_I\) is surjective; and for every \(I\)
  satisfying this equivalence, the fourth conjunct of the normal form
  holds if and only if the fifth does.
\end{enumerate}
\end{theorem}
\begin{proof}
We first show that a surjective \(\currQ_I\) gives
\(\operatorname{PredictiveSurplus}(I)\iff\neg\operatorname{PredictiveCollapse}(I)\);
this is \Cref{thm:F7} for the single probe \(\predM_I\), which we spell
out.  A surplus pair \(h,h'\) rules out \(\predM_I=g\circ\currQ_I\),
since it would give
\(\predM_I(h)=g(\currQ_I h)=g(\currQ_I h')=\predM_I(h')\).  Conversely,
if there is no surplus pair, \(\predM_I\) is constant on the fibers of
\(\currQ_I\), and since \(\currQ_I\) is surjective,
\(g(\currQ_I h):=\predM_I(h)\) defines \(g\) with
\(\predM_I=g\circ\currQ_I\).

(i) By the fourth conjunct, \(\neg\operatorname{PredictiveCollapse}(I)\)
is equivalent to
\(\neg\operatorname{lawfulCurrentExposure}(I)=\operatorname{HiddenFromCurrent}(I)\),
which with the equivalence just proved gives the fifth conjunct.  The
final chain combines the third and fourth conjuncts with the same
equivalence.

(ii) If \(\currQ_I\) is surjective, the equivalence was proved above.
If \(M_I\ne\emptyset\), the same argument applies, with the
fiber-defined map extended outside the image of \(\currQ_I\) by any
chosen value.  If \(M_I=\emptyset\), then \(\mathcal H_I=\emptyset\),
there is no surplus, and collapse holds exactly when the current
codomain is empty, that is, when \(\currQ_I\) is surjective.  Given the
equivalence, the fourth conjunct is equivalent to
\(\neg\operatorname{PredictiveSurplus}(I)\iff\operatorname{lawfulCurrentExposure}(I)\),
which is the fifth conjunct negated on both sides.
\end{proof}

\paragraph{Nonclaims.} The theorem does not prove \((\mathsf S)\) for
any interface derived from the saturated-SAT construction of
\citet{Tsiokos2026PvNP}, does not re-prove the closure result of that
paper, and asserts nothing about a complexity class.  Its unconditional
content is part (ii).  The upstream factorization here concerns the
whole declared readout; it is distinct from the exposure predicate of
the Hiddenness paper \citep{Tsiokos2026CSL}, which asks for a single
nonconstant visible function of state.  Since \(\predM_I\) refines
\(\currQ_I\), the pairing \(\langle\currQ_I,\predM_I\rangle\), the
coarsest refinement of \(\currQ_I\) through which \(\predM_I\) descends
(\Cref{sec:apparatus}), has the same fibers as \(\predM_I\).  The
Kochen--Specker theorem \citep{KochenSpecker1967Hidden}, a
non-contextual hidden-variable obstruction in quantum foundations, is
a substrate-specific result, not a special case of this law.

\subsection*{Layer instantiations}\label{layer-inst:F13a}

\begingroup\small
\begin{description}
\item[\emph{Hidden Markov model identifiability (statistical
latent-variable analysis).}] \textit{Analogy.}
In a hidden Markov model, two observation histories ending in the
same observed symbol can predict different distributions of future
observations, because they lead to different posterior distributions
over the hidden state; such a pair is a predictive surplus over the
current observation, and the hidden state is the determining state
upstream of it.  For a stationary hidden Markov model with \(r\) hidden
states and \(\kappa\ge2\) observation symbols, Allman, Matias and
Rhodes show that the parameters are generically identifiable, up to
relabeling of the hidden states, from the joint distribution of
\(2k+1\) consecutive observations whenever
\(\binom{k+\kappa-1}{\kappa-1}\ge r\); \emph{generically}
excludes a set of parameters of measure zero.  Identifiability concerns
the model, not the success of an estimation procedure such as
Baum--Welch EM \citep{AllmanMatiasRhodes2009Identifiability}.

\end{description}
\endgroup


\subsection{CSL Formed-Closure Law [\texorpdfstring{\Fref{F13b}}{F13b}]}\label{sec:access:csl-formed-closure}

\begin{definition}[Diagnosis source data]\label{def:diagnosis-source}
Diagnosis source data on a carrier \(H\) is a triple \((\currQ,\predM,d)\) of
a current quotient \(\currQ:H\to Q\), a predictive quotient
\(\predM:H\to M\) and a readout \(d:H\to S\) that is \emph{determining}:
\(\currQ(h)=\currQ(h')\wedge d(h)=d(h')\Rightarrow\predM(h)=\predM(h')\). In
the CSL case, \PaperCSL\ \citep{Tsiokos2026CSL} identifies an internal-state
readout; the determining property is an additional hypothesis that must be
verified for a proposed diagnosis source.
\end{definition}

\Cref{thm:F13b} applies to exactly such data.

\paragraph{Reading the statement.} The theorem has three parts: the surplus equivalence for the given determining \(d\); for an arbitrary \(d'\), a factorization test for being determining and a two-case split of predictive surplus; and the statement that \(\langle\currQ,\predM\rangle\) is the coarsest pairing \(\langle\currQ,d'\rangle\) with \(d'\) determining.
The closure carries a declared determining-state readout \(d\), the hidden
state that the diagnosis of \PaperCSL\ identifies.
\(\operatorname{PredictiveSurplus}(\mathcal C)\) says that two histories with
the same current class have different predictive classes;
\(\operatorname{HiddenStateSurplus}(\mathcal C)\) says that two histories with
the same current class and different determining states have different
predictive classes, so that the predictive split is carried by the hidden
determining state. The hypothesis that \(d\) is determining says that the
current class and the determining state together fix the predictive class.
This pairwise property is assumed here; it is not established by the cited CSL
diagnosis.

\begin{theorem}[CSL Formed-Closure Law]\label{thm:F13b}
Let \(\mathcal C=(H,\currQ,\predM,d)\) consist of a carrier \(H\) with
current quotient \(\currQ:H\to Q\), predictive quotient \(\predM:H\to M\) and a
declared determining-state readout \(d:H\to S\) that is determining:
\[
  \currQ(h)=\currQ(h')\ \wedge\ d(h)=d(h')
  \ \Longrightarrow\ \predM(h)=\predM(h').
\]
Define
\begin{align*}
  \operatorname{PredictiveSurplus}(\mathcal C)
  &:\Longleftrightarrow
  \exists h,h'\in H:\ \currQ(h)=\currQ(h')\ \wedge\ \predM(h)\ne\predM(h'),\\
  \operatorname{HiddenStateSurplus}(\mathcal C)
  &:\Longleftrightarrow
  \exists h,h'\in H:\ \currQ(h)=\currQ(h')\ \wedge\ d(h)\ne d(h')\\
  &\qquad{}\wedge\ \predM(h)\ne\predM(h').
\end{align*}
Then
\[
  \operatorname{PredictiveSurplus}(\mathcal C)
  \iff
  \operatorname{HiddenStateSurplus}(\mathcal C),
\]
and more precisely every pair \(h,h'\) with \(\currQ(h)=\currQ(h')\) and
\(\predM(h)\ne\predM(h')\) has \(d(h)\ne d(h')\): predictive surplus across
the \(\predM/\currQ\) gap lies only on pairs that the hidden determining state
separates; equivalently, the surplus pairs are the current-fiber pairs that both \(d\) and
\(\predM\) separate. Moreover, for an arbitrary readout \(d':H\to S'\), not assumed determining, \(d'\) is determining if and only if
\(\predM=g\circ\langle\currQ,d'\rangle\) for some map \(g\) from the image of
\(\langle\currQ,d'\rangle\) to \(M\) (equivalently, when \(M\) is nonempty, for some
\(g:Q\times S'\to M\)), and
\(\operatorname{PredictiveSurplus}\) holds if and only if some pair of one current
fiber is separated by both \(d'\) and \(\predM\), or some pair of one current
fiber agrees under \(d'\) and is separated by \(\predM\) (this second disjunct
says exactly that \(d'\) is not determining). In particular
\(d'=\predM\) is determining, and \(d'\) is determining exactly when
\(\langle\currQ,\predM\rangle\) is constant on the fibers of
\(\langle\currQ,d'\rangle\); hence, after corestricting the pairings to their images,
\(\langle\currQ,\predM\rangle\) is the coarsest pairing \(\langle\currQ,d'\rangle\) with
\(d'\) determining. When \(M\) is nonempty, the canonical pairing also factors
through the full-codomain pairing \(\langle\currQ,d'\rangle:H\to Q\times S'\) for
every determining \(d'\).
\end{theorem}
\begin{proof}
If \(\currQ(h)=\currQ(h')\) and \(\predM(h)\ne\predM(h')\), then
\(d(h)\ne d(h')\), since otherwise the determining hypothesis would give
\(\predM(h)=\predM(h')\). So every surplus pair is a hidden pair, which gives
\(\operatorname{PredictiveSurplus}\Rightarrow\operatorname{HiddenStateSurplus}\);
the converse holds because a hidden pair is in particular a surplus pair.
For an arbitrary \(d'\), if \(d'\) is determining, define \(g(a,b)\) as
\(\predM(h)\) for any \(h\) with \(\langle\currQ,d'\rangle(h)=(a,b)\), which is
well defined by the determining property; when \(M\) is nonempty, \(g\) extends
to \(Q\times S'\) by a fixed value off the image. Conversely such a \(g\) makes
\(d'\) determining. The surplus disjunction splits a surplus pair by
whether \(d'\) separates it.
For the last clause, \(\predM\) is determining trivially, and since the \(\currQ\) coordinates of two histories with equal \(\langle\currQ,d'\rangle\) agree, determination of \(d'\) is constancy of \(\langle\currQ,\predM\rangle\) on the fibers of \(\langle\currQ,d'\rangle\).
\end{proof}

\paragraph{What the source supplies.} \PaperCSL\ \citep{Tsiokos2026CSL} defines an internal-state readout and studies its exposure to current observables and its behavior under trace-fingerprint loops. It does not prove that \((\currQ,d)\) determines \(\predM\) for every pair of histories. That determination is an explicit hypothesis of \Cref{thm:F13b}; applying the theorem to a CSL carrier requires a separate proof of it.


\paragraph{Nonclaims.} This subsection does not instantiate the CSL diagnosis source and does
not assert a measurement-theoretic contextuality theorem at the
\FIV\ layer.

\subsection*{Layer instantiations}\label{layer-inst:F13b}

\begingroup\small
\begin{description}
\item[\emph{Bayesian clinical diagnostic testing (medicine).}] \textit{Analogy.}
For a clinical assay with declared sensitivity \(\mathrm{Se}\),
specificity \(\mathrm{Sp}\) and cohort prevalence, Bayes' rule gives
\(P(\text{disease}\mid r)\) for each possible result \(r\). Take the current
class to record the cohort and assay parameters, and the determining readout
to record \(r\); together they determine the posterior. If two possible
results give different posteriors within one current class, they witness
predictive and hidden-state surplus. Disease status alone need not determine
the posterior, and this analogy does not infer disease status from that
surplus \citep{Sackett2000EvidenceBasedMedicine}.

\end{description}
\endgroup


\subsection{Layer Multiplicity [\texorpdfstring{\Fref{F24}}{F24}]}\label{sec:access:layer-multiplicity}

Let \(H\) be a closure carrier with a formed current access quotient
\(q:H\to Q\).  A calibration or role signature is another declared
map \(s:H\to S\): it may record a future probe, an interaction role,
a lawhood signature, a branch readout, or a calibrated instrument
response.  The signature is already carried by the old layer exactly
when it descends through \(q\), equivalently when it is constant on
the \(q\)-fibers.

\paragraph{Reading the statement.}
\(\operatorname{RoleSplit}(q,s)\) says that the old layer merges two
histories whose roles differ, that is, some \(q\)-fiber contains two
values of \(s\). \(\operatorname{StrictLayerExtension}(q,q_s)\) says
that \(q_s\) retains \(q\), that is \(q=\mathrm{pr}_Q\circ q_s\), and that some
\(q\)-fiber is split by \(q_s\): there are \(h,h'\) with \(q(h)=q(h')\) and
\(q_s(h)\ne q_s(h')\). The theorem shows that a role split
is always resolved by this extension. The layer-formation discipline also
records eight other ways a split may be handled: by a memory layer, by a
hidden upstream role, through a bridge, within a budget, within a declared
scope, by coarsening, by placing the role outside the declared scope, or as
blocked nonclosure. These are declared statuses, not determined by \(q\) and
\(s\), and the theorem makes no claim about them. In ordinary notation: \(s\) factors
through \(q\) exactly when \(\mathcal O_s=\emptyset\); some \(q\)-fiber
carries two values of \(s\) exactly when \(\mathcal O_s\neq\emptyset\),
which happens exactly when \(q_s\) strictly refines \(q\); and a split
is then resolved at least by the strict extension \(q_s\). With \(d=s\), the
descent criterion, the strict refinement under a split and the coarsest-map
property are those of \Cref{thm:F12}; the converse (a strict extension
implies a role split) and the corestriction are added here.

\begin{theorem}[Layer Multiplicity]\label{thm:F24}
Let \(H\) be a history space, let \(q:H\to Q\) be a surjective access
quotient, and let \(s:H\to S\) be a declared role readout.  Define
the role obstruction and joint role-layer quotient by
\[
  \mathcal O_s
  :=
  \{(h,h')\in H^2:q(h)=q(h')\ \text{and}\ s(h)\ne s(h')\},
  \qquad
  q_s(h):=(q(h),s(h)).
\]
Then
\begin{align*}
  s\ \text{descends through}\ q
  &\iff \mathcal O_s=\emptyset,\\
  \operatorname{RoleSplit}(q,s)
  &\iff \mathcal O_s\ne\emptyset\\
  &\iff \operatorname{StrictLayerExtension}(q,q_s).
\end{align*}
In particular,
\[
  \operatorname{RoleSplit}(q,s)\Longrightarrow
  \operatorname{StrictLayerExtension}(q,q_s):
\]
a role split is resolved by the strict layer extension \(q_s\). Moreover, \(q_s:H\to Q\times S\) is coarsest in the factorization
preorder among maps through which both \(q\) and \(s\) factor: if
\(q=a\circ p\) and \(s=b\circ p\) for a map \(p\), then
\(q_s=\langle a,b\rangle\circ p\). Its corestriction
\(H\to\operatorname{im}(q_s)\) is the corresponding coarsest surjective
quotient among surjective maps \(p\) through which \(q\) and \(s\) factor.
\end{theorem}
\begin{proof}
Assume \(q:H\to Q\) is surjective and \(s:H\to S\) is declared.  If
\(s\) descends through \(q\), then \(s=\bar s\circ q\) for some
\(\bar s:Q\to S\), so
\[
  q(h)=q(h')\Longrightarrow s(h)=s(h')
\]
and \(\mathcal O_s=\emptyset\).  Conversely, if
\(\mathcal O_s=\emptyset\), define
\[
  \bar s(q(h)):=s(h).
\]
The empty obstruction makes \(\bar s\) independent of the
representative \(h\), and surjectivity of \(q\) supplies all quotient
classes.  Thus \(s\) descends through \(q\).

If \(\mathcal O_s\ne\emptyset\), choose
\((h,h')\in\mathcal O_s\).  Then \(q(h)=q(h')\) but
\[
  q_s(h)=(q(h),s(h))\ne(q(h'),s(h'))=q_s(h'),
\]
so \(q_s\) strictly refines \(q\).  Conversely, a strict refinement of
\(q\) by \(q_s\) means that some \(q\)-fiber is split by the second
coordinate \(s\), giving an element of \(\mathcal O_s\).

Finally, if \(\operatorname{RoleSplit}(q,s)\) holds, the previous
paragraph shows that \(q_s\) strictly refines \(q\), which is
\(\operatorname{StrictLayerExtension}(q,q_s)\).
For the last clause, \(q_s(h)=(q(h),s(h))=(a(p(h)),b(p(h)))=\langle a,b\rangle(p(h))\).
Corestricting \(q_s\) to its image makes it surjective. If \(p\) is
surjective, every \(t\) is \(p(h)\) for some \(h\), so
\((a(t),b(t))=q_s(h)\) lies in that image and gives the asserted
quotient factorization.
\end{proof}

\paragraph{Nonclaims.} The theorem proves that a role split forces
the strict extension \(q_s\); adopting \(q_s\) as a new layer still
requires admissibility, closure formation, and status.
Layer Multiplicity does not classify the number of layers in a
carrier, does not prescribe how layers must be presented, and does
not assert that multiplicity is finite in every setting.

\subsection*{Layer instantiations}\label{layer-inst:F24}

\begingroup\small
\begin{description}
\item[\emph{Operating-system privilege protection (computer
security).}] \textit{Analogy.}
In a multi-user operating system, the privilege quotient classifies
operations by the trust level of the executing principal.  When a
process needs an operation outside its declared privilege, the
layer-formation discipline names a resolution: a user-to-kernel
system-call trap, a capability-bridge invocation, an
inter-process-communication transfer, a sandbox or namespace
restriction, or
outright denial --- corresponding as follows to the declared statuses memory layer (the trap's
saved context), bridge (capability invocation and IPC transfer), declared scope
(sandbox) and blocked nonclosure (denial) listed in the reading paragraph of
\Fref{F24}, about which \Fref{F24} makes no claim
\citep{SaltzerSchroeder1975ProtectionInformation}.

\end{description}
\endgroup


\subsection*{Access at a glance}\label{sec:access:summary}\addcontentsline{toc}{subsection}{Access at a glance}
\begin{table}[H]
  \centering
  \footnotesize
  \begin{tabularx}{\textwidth}{@{}
      >{\raggedright\arraybackslash}p{0.24\textwidth}
      >{\raggedright\arraybackslash}X
      >{\raggedright\arraybackslash}p{0.12\textwidth}@{}}
    \toprule
    Name & One-line claim & Substrate cite \\
    \midrule
    Quotientality theorem (\Fref{F10}) &
    Every observable constant on access classes factors uniquely through the access quotient, and the access quotient factors through every surjection on whose fibers the signature is constant (it is the coarsest such quotient).  &
    \textemdash \\
    Adequacy / No-Overread theorem (\Fref{F11}) &
    Under the no-overread rule, accepted exact claims factor uniquely through the access quotient.  &
    \textemdash \\
    No-Free-Distinction theorem (\Fref{F12}) &
    A readout is current-visible iff it factors through the access
    quotient; the pairing with it is always the coarsest map through which
    both factor, strict when the readout splits a fiber, and then, under
    (NO), the readout is not accepted as exact.  &
    \textemdash \\
    Hiddenness Normal Form (\Fref{F13a}), conditional &
    Under the source hypothesis \((\mathsf S)\) (surjectivity and four of
    the five conjuncts), the fifth conjunct holds, and exposure, collapse
    and absence of surplus are equivalent; without \((\mathsf S)\),
    predictive surplus holds exactly when predictive collapse fails, for a
    surjective current quotient. &
    \PaperPvNP \\
    CSL Formed-Closure (\Fref{F13b}) &
    With a determining-state readout, every predictive-surplus pair is separated by the determining state.  &
    \PaperCSL \\
    Layer-Multiplicity Theorem (\Fref{F24}) &
    A role readout splits an access quotient exactly when its role
    obstruction is nonempty, and a role split implies the strict
    layer extension by the joint role-layer quotient.  &
    \textemdash \\
    \bottomrule
  \end{tabularx}
  \caption{Access at a glance: the six laws of \Cref{sec:access}.}
  \label{tab:axis-access-summary}
\end{table}

\FloatBarrier
